\section{An extension to symmetric matrices}
\label{sec:extension}

Fix an admissible map $f:\mathbb P(V)\to\mathbb R^m$, where
$\dim V=k$.  We extend $f$ first to positive semidefinite operators and
then to all symmetric operators.  Orthogonal spectral subspaces will give
disjoint sets of labels, producing an odd map between spheres.

Define the continuous degree-two homogeneous function
$H:V\to\mathbb R^m$ by
\[
 H(v)=\begin{cases}\|v\|^2f([v]),&v\ne0,\\0,&v=0.\end{cases}
\]
Continuity at zero follows from $\|H(v)\|_1=\|v\|^2$.
For a positive semidefinite operator $Q$, let
\begin{equation}
 E_V(Q)=k\int_{\mathbb S(V)}H(Q^{1/2}u)\,d\mu_V(u),
 \label{eq:extension}
\end{equation}
where $\mu_V$ is rotation-invariant probability measure on the unit sphere.
The positive square root is meant throughout.  We usually write $E$
when the ambient space is clear.

\begin{lemma}[Properties of the positive extension]
\label{lem:positive-extension}
The map $E_V$ is continuous on the full positive semidefinite cone,
including its singular boundary.  For $t\geq0$ it satisfies
\[
 E_V(tQ)=tE_V(Q),\qquad E_V(Q)_i\geq0,\qquad
 \sum_{i=1}^m E_V(Q)_i=\operatorname{tr}Q.
\]
Its positive coordinates are exactly the labels used on the range of $Q$:
\begin{equation}
 \{i:E_V(Q)_i>0\}
 =\bigcup_{x\in\mathbb P(\operatorname{ran}Q)}S(x),
 \label{eq:extension-support}
\end{equation}
where the right side is empty for $Q=0$.  Moreover:
\begin{enumerate}
\item $E_V(P_x)=f(x)$ for every line $x$;
\item if $Q$ is supported on a nonzero subspace $W\subseteq V$, then
$E_V(Q)=E_W(Q|_W)$, where the latter construction uses $f|_{\mathbb P(W)}$
and the factor $\dim W$;
\item this last expression is smooth when $W$ varies smoothly with fixed
dimension and $Q|_W$ varies smoothly through positive definite operators.
\end{enumerate}
\end{lemma}

\begin{proof}
The positive square root depends continuously on a positive semidefinite
matrix.  For example, approximate the square-root function uniformly by
polynomials on a compact interval containing the spectra under
consideration, and apply those polynomials to the matrices.
Thus the integrand in \eqref{eq:extension} is jointly continuous in
$(Q,u)$; integration over the compact sphere proves continuity of $E_V$.
Its homogeneity and nonnegativity follow from the corresponding properties
of $H$.  Since $\sum_i H(v)_i=\|v\|^2$ and
$\int uu^{\mathsf T}\,d\mu_V(u)=k^{-1}I$, we have
\[
 \sum_i E_V(Q)_i=k\int\langle Qu,u\rangle\,d\mu_V(u)
 =\operatorname{tr}Q.
\]

Every nonzero $Q^{1/2}u$ lies in $\operatorname{ran}Q$.
Conversely, every line in that range is spanned by some $Q^{1/2}u$ with
$\|u\|=1$, because $Q^{1/2}$ is invertible on its range.
If $f_i$ is positive on such a line, the $i$th integrand is positive at
that $u$, and hence on a nonempty open set of the sphere.  Such a set has
positive measure.  This proves \eqref{eq:extension-support} in both
directions.  For $Q=P_x$, every nonzero integrand is a scalar multiple
of $f(x)$, and the scalars integrate to $\operatorname{tr}P_x=1$.
Thus $E_V(P_x)=f(x)$.

To prove the subspace identity, write $r=\dim W$ and let $\pi_W$ be
orthogonal projection.  The finite measure on $\mathbb S(W)$ defined by
\[
 B\longmapsto\int_{\mathbb S(V)}
 \|\pi_Wu\|^2\,\mathbf1_{\{\pi_Wu\ne0,\;
                  \pi_Wu/\|\pi_Wu\|\in B\}}\,d\mu_V(u)
\]
is invariant under all orthogonal transformations of $W$.  Its total mass is
$\int\|\pi_Wu\|^2\,d\mu_V(u)=r/k$, so it equals
$(r/k)\mu_W$.  Because $Q$ vanishes on $W^\perp$ and $H$ is homogeneous
of degree two, this gives
\[
 k\int_{\mathbb S(V)}H(Q^{1/2}u)\,d\mu_V(u)
 =r\int_{\mathbb S(W)}H((Q|_W)^{1/2}w)\,d\mu_W(w),
\]
as claimed.

Finally choose, locally in the parameters, a smooth orthonormal frame
$T:\mathbb R^r\to W$.  Write $Q|_W$ in that frame as the positive
definite matrix $C$.  The subspace formula becomes
\[
 E_V(Q)=r\int_{\mathbb S^{r-1}}H(TC^{1/2}w)\,d\mu(w).
\]
The matrix square root is smooth on positive definite matrices, and
$TC^{1/2}w$ never vanishes.  The function $H$ is smooth away from zero.
Differentiating this expression over the compact sphere proves the
asserted local smoothness.
\end{proof}

For a symmetric operator $A$, let $|A|=(A^2)^{1/2}$ and
$A_+=(|A|+A)/2$, $A_-=(|A|-A)/2$.  These are its positive and negative
parts, so $A=A_+-A_-$ and their ranges are orthogonal.  Define
\begin{equation}
 F(A)=E_V(A_+)-E_V(A_-).
 \label{eq:signed-extension}
\end{equation}

\begin{lemma}[The odd sphere map]
\label{lem:signed-extension}
The map $F:\operatorname{Sym}(V)\to\mathbb R^m$ is continuous, odd,
and homogeneous of degree one for nonnegative scalars.  It satisfies
\begin{equation}
 \|F(A)\|_1=\operatorname{tr}|A|,
 \qquad \sum_iF(A)_i=\operatorname{tr}A.
 \label{eq:extension-norms}
\end{equation}
More precisely, its positive and negative label sets are
\begin{equation}
\begin{aligned}
 \{i:F(A)_i>0\}&=\bigcup_{x\in\mathbb P(\operatorname{ran}A_+)}S(x),\\
 \{i:F(A)_i<0\}&=\bigcup_{x\in\mathbb P(\operatorname{ran}A_-)}S(x).
\end{aligned}
 \label{eq:spectral-support}
\end{equation}
If every coordinate of $F(A)$ is nonzero, then $A$ is nonsingular.
If all coordinates of $F(A)$ are strictly negative, then $A$ is negative
definite; the analogous statement holds with both signs reversed.
The map $F$ is smooth on the open set of nonsingular symmetric operators.
In particular, it restricts to a continuous odd map
\begin{equation}
 \{A:\operatorname{tr}|A|=1\}\longrightarrow
 \{y\in\mathbb R^m:\|y\|_1=1\}.
 \label{eq:norm-sphere-map}
\end{equation}
\end{lemma}

\begin{proof}
Continuity and homogeneity follow from Lemma~\ref{lem:positive-extension}
and the continuous formulas for $A_\pm$.  Replacing $A$ by $-A$ exchanges
$A_+$ and $A_-$, so $F$ is odd.  Any line in
$\operatorname{ran}A_+$ is orthogonal to every line in
$\operatorname{ran}A_-$.  Admissibility and
\eqref{eq:extension-support} therefore imply that $E_V(A_+)$ and
$E_V(A_-)$ have disjoint positive coordinates.  There is no cancellation
in their difference.  This proves \eqref{eq:spectral-support} and, by
summing coordinates, \eqref{eq:extension-norms}.

If $\ker A$ contains a line $x$, choose $i\in S(x)$, which is possible
because $\sum_i f_i(x)=1$.  The line $x$ is orthogonal to both spectral
ranges.  Thus label $i$ is absent from both unions in
\eqref{eq:spectral-support}, so $F(A)_i=0$.  This proves the
nonsingularity assertion.  If all coordinates of $F(A)$ are negative,
then $E_V(A_+)=0$ by disjointness of the two positive supports.  Hence
$\operatorname{tr}A_+=0$, so $A_+=0$, and nonsingularity makes $A$
negative definite.  Oddness gives the positive case.

For local smoothness, fix nonsingular $A$.  In a neighborhood of $A$,
the positive and negative spectral subspaces have constant dimensions
and depend smoothly on the matrix.  Explicitly their projections are
\[
 \Pi_\pm(A)=\frac12\bigl(I\pm A(A^2)^{-1/2}\bigr),
\]
which are smooth there because $A^2$ is positive definite.
Each projection admits a local smooth orthonormal frame: project a fixed
basis for its range and apply Gram--Schmidt after shrinking the
neighborhood.  In these frames, $A_+$ and $A_-$ restrict to positive
definite operators on their respective ranges.  The last assertion of
Lemma~\ref{lem:positive-extension} proves smoothness of $E_V(A_+)$ and
$E_V(A_-)$, with a zero term if the corresponding range has dimension
zero.  This proves smoothness of $F$.
\end{proof}

The two norm spheres in \eqref{eq:norm-sphere-map} are not being regarded
as globally smooth hypersurfaces.  Radial normalization identifies them
antipodally with the ordinary Euclidean spheres of dimensions $N_k-1$
and $m-1$.  These identifications are homeomorphisms everywhere and
smooth diffeomorphisms locally at nonsingular matrices and label vectors
with every coordinate nonzero.  Indeed, $\operatorname{tr}|A|$ is smooth
near a nonsingular $A$, while the $\ell^1$ norm is a fixed linear form
within each open sign orthant.  Consequently ordinary sphere degree
applies to \eqref{eq:norm-sphere-map} when $m=N_k$, and local differential arguments
are valid at every preimage of a full-support value.
